\documentclass[11pt]{article}
\usepackage[margin=0.8in]{geometry}
\usepackage{amsmath,amssymb,amsthm}
\usepackage{booktabs}
\usepackage{enumitem}
\usepackage{xcolor}
\usepackage{hyperref}
\hypersetup{colorlinks=true,linkcolor=blue!55!black}

\title{\textbf{JSET v1} — One-Page Summary\\
\large Stone–Jordan Exceptional Theory (SJET$_{\mathrm{true}}$): Honest Ledger for an Exceptional Jordan/E$_6$ Scaffold}
\author{Brandon Belna}
\date{June 2026}

\begin{document}
\maketitle

\noindent\textbf{What JSET v1 is (and is not).} JSET v1 is a \emph{conditional} framework for organizing GUT-like matter, family/observer geometry, and effective dynamics using the exceptional Jordan algebra and $E_6$. It is \textbf{not} a derivation of the Standard Model, general relativity, QFT, or cosmology from two primitives (void + mirror). All stronger claims from earlier versions have been restructured following a 39-agent adversarial validation (see \texttt{SJET-VALIDATION-REPORT.md}).

\section*{The Rigorous Core (Standard Imported Theorems)}

\begin{itemize}[leftmargin=*]
\item Albert algebra $J_3(\mathbb{O})$: $\dim_{\mathbb{R}} J = 3 + 3 \cdot 8 = 27$.
\item EIII Hermitian symmetric space: $X = E_6^{\mathrm{cpt}} / (\Spin(10) \cdot U(1))$, $\dim_{\mathbb{R}} X = 78 - 45 - 1 = 32$.
\item Tangent representation: $T_o X \otimes_{\mathbb{R}} \mathbb{C} \cong \mathbf{16}_{-3} \oplus \overline{\mathbf{16}}_{+3}$.
\item Minuscule branching (standard):
  \[
  \mathbf{27} \downarrow_{\Spin(10)\times U(1)} = \mathbf{1}_{+4} \oplus \mathbf{10}_{-2} \oplus \mathbf{16}_{+1}.
  \]
\end{itemize}

These are classical facts from Jordan algebra and Lie theory. They carry no physics by themselves.

\section*{Typed Mirror Scaffold (No Single Functor)}

The operations appearing in the exceptional ladder (Boolean doubling, Cayley–Dickson, Jordan exit, Freudenthal–Tits, product, EIII reduction) are \textbf{qualitatively different}. There is currently no single involutive functor $\Mirror$ with $\Mirror^2 \simeq \id$ that uniformly generates the ladder. 

\textbf{Corrected claim:} JSET studies a \emph{typed} exceptional-algebra scaffold. The existence of a unified categorical mirror is an open conjecture.

\section*{Capacity Layer (Bookkeeping, Not Derivation)}

The capacity functional
\[
\mathcal{C}(\omega) = -\sum_a \mu(a) \log\Bigl(\sum_i e^{c_i(a)-\omega_i}\Bigr), \quad \sum_i \omega_i = 0
\]
is strictly concave on the hyperplane under mild assumptions. Its constrained maximizer on uniform data yields the uniform distribution $1/N$. The specific weights $(1/27, 10/27, 16/27)$ are the normalized dimensions of the $E_6 \to \Spin(10)\times U(1)$ branching fed back into the functional (or supplied via a Legendre term $\mathcal{L}_m$). 

\textbf{Corrected claim:} Capacity is a convex bookkeeping (or moment-map) device for an \emph{already chosen} sector split unless the graph, measure $\mu$, and costs $c_i$ are specified independently of the target irrep dimensions.

\section*{Family/Observer Geometry — Conjectural Axiom \textbf{FAM}}

To have a well-defined family count $n_F = \dim H^1(Q, V_{16})^{-}$ one needs: a compact complex/algebraic space $Q$, a genuine involution $\sigma$ with $\sigma^2 = \id$ (order exactly 2), a $\sigma$-equivariant holomorphic bundle $V_{16}$, and an explicit computation of the anti-invariant cohomology.

\textbf{\textbf{FAM} (Conjecture):} There exists $(Q, \sigma, V_{16})$ satisfying the above with $\dim H^1(Q, V_{16})^{-} = 3$.

Previous claims that this dimension equals 3 were not supported by a correct computation (incoherent cellular quotient, order-6 vs order-2 inconsistency for $\sigma$, misapplied Lefschetz formula with no fixed points implying index 0).

\section*{Spacetime Soldering — Corrected}

For normed division algebra $\mathbb{K}$:
\[
\dim_{\mathbb{R}} H_2(\mathbb{K}) = 2 + \dim_{\mathbb{R}} \mathbb{K}, \qquad \text{signature of }\det = (1, 1 + \dim_{\mathbb{R}}\mathbb{K}).
\]

Explicitly: $H_2(\mathbb{C})$ gives 4d Minkowski (signature $(1,3)$); $H_2(\mathbb{H})$ is 6-dimensional of signature $(1,5)$; $H_2(\mathbb{O})$ is 10-dimensional of signature $(1,9)$.

\textbf{Corrected claim:} Four-dimensional Lorentzian spacetime requires an explicit selection of a complex subalgebra $\mathbb{C} \hookrightarrow \mathbb{O}$ (or equivalent 4-plane) inside the Jordan structure. This selection is additional structure (named Assumption A-Solder-C), not derived by "quaternionic soldering."

\section*{Conditional Layers (EFT Only)}

Under named assumptions (E6 embedding + \textbf{FAM} for matter; OS axioms verified on the collapsed measure for QFT; explicit beta-function coefficients for gravity, etc.), one obtains:
- Standard representation-theoretic organization of three $\Spin(10)$ $\mathbf{16}$s (with correct tensor products $\mathbf{16}\otimes\mathbf{16} = \mathbf{10}_s \oplus \mathbf{120}_a \oplus \mathbf{126}_s$).
- Osterwalder–Schrader reconstruction \emph{if} the axioms hold for the concrete measure.
- Einstein–Hilbert truncations whose fixed-point values follow algebraically from assumed anomalous-dimension inputs.

None of these are forced by the exceptional core alone. No derivation of the full SM Lagrangian, Yukawa matrices, CKM/PMNS values, or quantum gravity is claimed.

\section*{The Six Open Problems (JSET v2 Roadmap)}

\begin{enumerate}[leftmargin=*]
\item Unified mirror category (single involutive functor recovering the typed operations with $\Mirror^2\simeq\mathrm{id}$ where needed).
\item Independent capacity data (graph $G$, $\mu$, $c_i$ that yield the 1|10|16 weights without circular input from the $E_6$ branching).
\item Explicit $(Q,\sigma,V_{16})$ + rigorous computation of $\dim H^1(Q,V_{16})^{-}$.
\item Mechanism (or postulate) selecting $\mathbb{C}\subset\mathbb{O}$ for correct 4d Lorentzian soldering.
\item Derivation (or justification) of the specific EFT operators and RG coefficients from the core + geometric assumptions.
\item Whether the framework yields any \emph{new} quantitative constraints beyond ordinary $E_6$/$\Spin(10)$ GUT model-building with comparable parameters.
\end{enumerate}

Progress on any of these can be measured precisely. Stronger physical claims require discharging (or openly retaining) the corresponding assumptions.

\section*{How to Read JSET v1}

1. Read the abstract and Section 3 (Rigorous Core) of the full ledger \texttt{sjet\_true.tex}.
2. Read the six open problems (Section 10).
3. Consult the validation report for the defects that necessitated the restructure.
4. Treat every physical statement as conditional on the explicitly named assumptions listed in the ledger.

\textbf{Citation note.} Cite \texttt{sjet\_true.tex} (or its PDF) for the rigorous core and the precise statement of the open problems. Do not cite it as evidence of a derivation of four-dimensional physics from void and mirror.

\vspace{0.5em}
\noindent\textit{This one-pager is extracted from the JSET v1 canonical document \texttt{sjet\_true.tex}. For full definitions, proofs of the core facts (as SITs), named assumptions, and the complete ledger classification, see the main file.}
\end{document}